% ============================================================
% 05-results.tex — V13 Results (condensed)
% ============================================================
\section{Results}
\label{sec:results}

%------------------------------------------------
\subsection{Scenario 1 --- Architecture-Level Comparison}
\label{sec:s1_results}

Table~\ref{tab:arch_comparison} compares the three configurations on the 100-query held-out benchmark. All share the same LLM, knowledge sources, and retrieval stack; differences are restricted to orchestration.

\begin{table}[!htbp]
\centering
\caption{Scenario~1: Architecture-level comparison ($N=100$). Tools/Q = mean tool invocations ($\downarrow$). Bold = best reliability / lowest cost.}
\label{tab:arch_comparison}
{\small\setlength{\tabcolsep}{3.5pt}
\begin{tabular}{lcccccccc}
\toprule
\textbf{Arch.} & \textbf{E2E} & \textbf{Fact} & \textbf{Src.R} & \textbf{Src.AP} & \textbf{Tools/Q} & \textbf{In Tok} & \textbf{Out Tok} & \textbf{Lat.\,(ms)} \\
\midrule
Single Agent   & 73.0 & 63.8 & 78.6 & 70.8 & 0.41 & 6296 & 1530 & \textbf{5159} \\
Routed Gen.    & 71.0 & 59.6 & 78.6 & 70.8 & 0.31 & \textbf{3473} & 1415 & 8921 \\
\textbf{CTU-Chat} & \textbf{81.0} & \textbf{66.7} & 78.6 & 70.8 & \textbf{0.28} & 3690 & \textbf{2092} & 10652 \\
\bottomrule
\end{tabular}}
\end{table}

\textbf{Pairwise bootstrap analysis} (10,000 resamples):
(1)~CTU-Chat vs.\ Single Agent: CTU-Chat achieves an 8.0\,pp higher E2E point estimate (81.0\% vs.\ 73.0\%, paired 95\%\,bootstrap CI $[0.0, +16.0]$\,pp; exact McNemar $p{=}0.096$), alongside higher Fact Coverage ($\Delta{=}+0.028$, 95\%\,CI $[-0.021,+0.079]$) and a 41.4\% reduction in input tokens ($\Delta{=}-2607$, CI $[-3325,-1868]$), at the cost of $+5.49$\,s median end-to-end latency.
(2)~CTU-Chat vs.\ Routed Generic: specialist policies raise Fact Coverage by $+7.05$\,pp (CI $[+0.010,+0.132]$, $p{<}0.05$).
(3)~Routed Generic vs.\ Single Agent: comparable token savings ($\Delta{=}-2824$) but lower fact coverage ($\Delta{=}-0.042$, CI $[-0.098,+0.013]$), indicating that routing without policy specialization is insufficient.

On the 20 cross-domain queries, CTU-Chat achieves an 85.0\% E2E point estimate compared with 75.0\% (Single) and 70.0\% (Generic), consistent with the intended single-owner multi-evidence-prefetch design.

%------------------------------------------------
\subsection{Scenario 2 --- Architectural Ablation}
\label{sec:s2_results}

Table~\ref{tab:ablation} removes one mechanism at a time from Full CTU-Chat ($N{=}100$). Deltas are relative to the full system.

\begin{table}[!htbp]
\centering
\caption{Scenario~2: Ablation results ($N=100$). Negative $\Delta$ = degradation.}
\label{tab:ablation}
{\small\setlength{\tabcolsep}{3pt}
\begin{tabular}{lccccc}
\toprule
\textbf{Ablated} & \textbf{$\Delta$E2E} & \textbf{$\Delta$Fact} & \textbf{$\Delta$Src.R} & \textbf{$\Delta$Tok} & \textbf{$\Delta$Lat} \\
\midrule
No Spec.~Policy  & $-$10.0 & $-$7.05 & 0.0  & $-$217  & $-$1731 \\
Full Tool Vis.   & $-$3.0  & $-$0.02 & 0.0  & +1270   & +296 \\
Uniform Text-RAG & $-$10.0 & $-$2.67 & $-$12.3 & $-$300  & $-$1910 \\
No Route Rep.    & $-$4.0  & $-$2.05 & 0.0  & $-$219  & $-$1610 \\
No Determ.~Calc.$^\dagger$ & $-$6.0 & $-$0.97 & 0.0 & $-$782 & $-$4540 \\
\bottomrule
\end{tabular}}
\smallskip
{\footnotesize $^\dagger$Full benchmark with calculators disabled. $\Delta$Tok/Lat in tokens/ms per query.}
\end{table}

The two most impactful ablations are \emph{No Specialist Policy} ($-10.0$\,pp E2E, $-7.05$\,pp Fact) and \emph{Uniform Text-RAG Baseline} ($-10.0$\,pp E2E, $-12.3$\,pp Source Recall). Replacing representation-specific access with uniform text retrieval substantially reduces source recall in the evaluated benchmark, disproportionately harming structured-financial queries ($-17.8$\,pp) and cross-domain queries ($-20.0$\,pp), while No Route Repair degrades cross-domain queries by $-15.0$\,pp. Disabling deterministic calculation produces a $-6.0$\,pp secondary drop across the full benchmark; on the arithmetic tuition subset ($N{=}9$), failure rate rises from 11.1\% to 66.7\% ($-55.6$\,pp exact match pass) due to unconstrained arithmetic hallucination.

\begin{table}[!htbp]
\centering
\caption{Scenario~2: E2E Success by query family across ablations.}
\label{tab:ablation_breakdown}
{\small\setlength{\tabcolsep}{4pt}
\begin{tabular}{lcccc}
\toprule
\textbf{Config.} & \textbf{Graph} & \textbf{Struct-fin} & \textbf{Narr-reg} & \textbf{Cross} \\
 & ($n{=}19$) & ($n{=}28$) & ($n{=}33$) & ($n{=}20$) \\
\midrule
\textbf{Full CTU-Chat} & 73.7 & \textbf{85.7} & \textbf{78.8} & \textbf{85.0} \\
No Spec.~Policy  & 63.2 & 82.1 & 66.7 & 70.0 \\
Full Tool Vis.   & \textbf{78.9} & 71.4 & \textbf{78.8} & \textbf{85.0} \\
Uniform Text-RAG & 73.7 & 67.9 & 75.8 & 65.0 \\
No Route Rep.    & \textbf{78.9} & 78.6 & 75.8 & 70.0 \\
\bottomrule
\end{tabular}}
\end{table}

%------------------------------------------------
\subsection{Scenario 3 --- Tool Ambiguity Stress Test}
\label{sec:s3_results}

Table~\ref{tab:tool_ambiguity} reports the semantic-neighbor experiment across 11--51 tools ($N{=}50$, $R{=}3$ order blocks; 1,800 primary $+$ 480 coverage decisions).

\begin{table}[!htbp]
\centering
\caption{Scenario~3: Semantic-neighbor tool ambiguity ($N=50$, $R=3$).}
\label{tab:tool_ambiguity}
{\small\setlength{\tabcolsep}{3pt}
\begin{tabular}{llccccc}
\toprule
\textbf{Reg.} & \textbf{Config.} & \textbf{Vis.} & \textbf{Sel.} & \textbf{E2E} & \textbf{Tok} & \textbf{Lat} \\
\midrule
\multirow{4}{*}{11} & Full-reg single & 11 & \textbf{94.0} & \textbf{86.0} & 1663 & \textbf{1098} \\
 & Top-10 single   & 10 & 93.3 & 83.3 & 1541 & 1148 \\
 & Routed gen.     & 5.1 & 90.7 & 80.0 & \textbf{1204} & 1889 \\
 & \textbf{Routed spec.} & 5.1 & 93.3 & 85.3 & 1441 & 1983 \\
\midrule
\multirow{4}{*}{31} & Full-reg single & 31 & 91.3 & 84.7 & 2706 & 1157 \\
 & Top-10 single   & 10 & 90.0 & 81.3 & \textbf{1149} & \textbf{1033} \\
 & Routed gen.     & 9.8 & 88.0 & 79.3 & 1396 & 1958 \\
 & \textbf{Routed spec.} & 9.8 & \textbf{93.3} & \textbf{88.7} & 1830 & 2015 \\
\midrule
\multirow{4}{*}{51} & Full-reg single & 51 & 92.7 & 83.3 & 3722 & 1289 \\
 & Top-10 single   & 10 & 86.7 & 78.7 & \textbf{1036} & \textbf{997} \\
 & Routed gen.     & 9.8 & 91.3 & 82.7 & 1340 & 1907 \\
 & \textbf{Routed spec.} & 9.8 & \textbf{92.7} & \textbf{88.0} & 1770 & 2021 \\
\bottomrule
\end{tabular}}
\end{table}

As the registry scales from 11 to 51 tools, the Full-Registry Single Agent degrades from 86.0\% to 83.3\% E2E while input tokens expand by 124\%. The Global Top-10 drops to 78.7\% (coinciding with shortlist truncation). The Routed Specialist preserves 88.0\% E2E at 51 tools with only 9.8 visible tools, maintaining a fixed-registry endpoint advantage (+4.7\,pp over Single Agent, +9.3\,pp over Top-10) without displaying statistically divergent degradation slopes.

%------------------------------------------------
\subsection{Supporting Retrieval Validation}
\label{sec:retrieval_results}

Table~\ref{tab:retrieval_evaluation} validates the retrieval subsystem independently.

\begin{table}[!htbp]
\centering
\caption{Supporting evaluation: Progressive retrieval component stacking ($N=100$).}
\label{tab:retrieval_evaluation}
{\small\setlength{\tabcolsep}{3pt}
\begin{tabular}{lcccccc}
\toprule
\textbf{Config.} & \textbf{Hit@1} & \textbf{Hit@3} & \textbf{P@5} & \textbf{R@5} & \textbf{MRR@10} & \textbf{Lat.\,(ms)} \\
\midrule
BM25 Lexical          & 0.810 & 0.920 & \textbf{0.234} & 0.937 & 0.874 & \textbf{14} \\
Dense Semantic        & 0.630 & 0.840 & 0.202 & 0.853 & 0.736 & 80 \\
Hybrid RRF            & 0.770 & 0.940 & 0.222 & 0.935 & 0.851 & 80 \\
Hybrid + Reranker     & 0.830 & 0.940 & 0.230 & 0.955 & 0.896 & 1{,}145 \\
\textbf{Full Retrieval} & \textbf{0.850} & \textbf{0.960} & \textbf{0.234} & \textbf{0.975} & \textbf{0.913} & 1{,}153 \\
\bottomrule
\end{tabular}}
\end{table}

The progressive retrieval ablation confirms that each component contributes complementary strengths. BM25 provides a competitive baseline (0.810 Hit@1, 0.874 MRR@10) due to the frequency of rigid identifiers in Vietnamese administrative regulations (decree numbers, course codes). Dense semantic search exhibits lower recall (0.853 R@5) from vocabulary mismatch on specialized institutional nomenclature. Hybrid RRF fusion expands candidate breadth (0.940 Hit@3), though top-rank precision is slightly diluted by dense candidates relative to pure BM25. Incorporating cross-encoder neural reranking (\texttt{bge-reranker-v2-m3}) overcomes this dilution, surpassing the BM25 baseline (0.830 Hit@1, 0.896 MRR@10). Finally, graph grounding further lifts full-stack retrieval performance to 0.850 Hit@1 and 0.975 Recall@5, indicating that structured knowledge constraints effectively resolve institutional entity ambiguities.

%------------------------------------------------
\subsection{Failure Analysis}
\label{sec:failure_analysis}

Table~\ref{tab:failures} categorizes error modes across the 46 failure instances ($<50\%$ Fact Coverage) in Scenario~1 (27 for Single Agent, 19 for CTU-Chat).

\begin{table}[!htbp]
\centering
\caption{Failure taxonomy across architecture configurations ($N=100$). Categories are non-mutually exclusive; a single query may exhibit multiple error modes.}
\label{tab:failures}
{\small
\begin{tabular}{llcc}
\toprule
\textbf{Category} & \textbf{Failure Type} & \textbf{Single} & \textbf{CTU} \\
\midrule
\multirow{2}{*}{Routing} & Misroute corrected by repair & N/A & 4 \\
 & Unrecoverable misroute & N/A & 2 \\
\midrule
\multirow{3}{*}{Retrieval} & Gold source absent & 2 & 2 \\
 & Cross-domain partial omission & 5 & 3 \\
 & Reranker demoted gold below $k{=}5$ & 4 & 4 \\
\midrule
\multirow{3}{*}{Tool} & Invalid/default arguments & 4 & 1 \\
 & Wrong-family hallucination & 3 & 1 \\
 & Unnecessary invocation & 5 & 2 \\
\midrule
\multirow{2}{*}{Generation} & Tool result ignored in text & 2 & 1 \\
 & Partial fact omission & 11 & 9 \\
\midrule
\textbf{Total} & Fact Coverage $<50\%$ & \textbf{27} & \textbf{19} \\
\bottomrule
\end{tabular}}
\end{table}

CTU-Chat's advantage is concentrated in tool-execution reliability (2 tool errors vs.\ 7; 2 unnecessary calls vs.\ 5). Both architectures share identical retrieval omissions (6 cases), confirming that retrieval quality is an orthogonal bottleneck.

